\documentclass[11pt,a4paper]{amsart}
\pdfinfoomitdate=1
\pdftrailerid{}
\pdfsuppressptexinfo=15
\usepackage[T1]{fontenc}
\usepackage{lmodern}
\usepackage[margin=27mm]{geometry}
\usepackage{microtype}
\usepackage{amsmath,amssymb,mathtools}
\usepackage{booktabs}
\usepackage{xcolor}
\usepackage[numbers,sort&compress]{natbib}
\usepackage[colorlinks=true,linkcolor=blue!55!black,citecolor=blue!55!black,urlcolor=blue!55!black]{hyperref}
\usepackage[nameinlink,noabbrev]{cleveref}
\raggedbottom

\newtheorem{theorem}{Theorem}[section]
\newtheorem{proposition}[theorem]{Proposition}
\newtheorem{corollary}[theorem]{Corollary}
\newtheorem{lemma}[theorem]{Lemma}
\theoremstyle{definition}
\newtheorem{definition}[theorem]{Definition}
\theoremstyle{remark}
\newtheorem{remark}[theorem]{Remark}

\newcommand{\T}{\mathbb T}
\newcommand{\N}{\mathbb N}
\newcommand{\Fix}{\operatorname{Fix}}
\newcommand{\card}[1]{\lvert #1\rvert}
\newcommand{\mob}{\mu}
\newcommand{\AM}{\mathrm{AM}}
\newcommand{\Lef}{\mathrm{L}}

\title[Finite-subset circle expansion]{Finite-Subset Circle Expansion:\
Parity-Split Fixed Counts and Alternating Zeta Factors}
\author{Anonymous}
\date{Internal Stage 2 draft, 28 August 2026}
\subjclass[2020]{37C30, 37B40, 54B20, 05A15}
\keywords{finite subset space, circle endomorphism, Artin--Mazur zeta function,
periodic orbit, configuration space, topological entropy}

\begin{document}

\begin{abstract}
Let $m_d(x)=dx\pmod 1$ on the circle, where $d\geq2$, and let
$H_{d,k}$ send a nonempty subset of at most $k$ circle points to its image
under $m_d$.  This gives a nonlinear, cardinality-collapsing map on the
$k$-dimensional finite-subset space.  We compute its complete unsigned
periodic ledger.  If $Q=d^n$, then fixed subsets of exact cardinality $j$
number
\[
 E_j(Q)=\frac{(Q-1)(Q^j-(-1)^j)}{Q+1}.
\]
Summing through $j=k$ yields $Q(Q^k-1)/(Q+1)$ for even $k$ and
$(Q^{k+1}-1)/(Q+1)$ for odd $k$.  The resulting Artin--Mazur zeta function
is a finite alternating product of the factors $1-d^r z$.  We also obtain
an exact M\"obius census and a prime-orbit asymptotic, prove
$h_{\rm top}(H_{d,k})=k\log d$ through a uniformly finite-to-one product
cover, and recover $(d,k)$ from the zeta function.  The proof rests on a
binary Euler transform: a fixed finite subset is a disjoint union of base
cycles, so each base cycle is either absent or present once.  Independent
coefficient and literal rational-circle enumerations verify the registered
finite controls.
\end{abstract}

\maketitle

\section{Introduction}

Finite-subset dynamics retains the motion of individual points while adding
collisions: distinct points may acquire the same image, and the cardinality
of a configuration may drop.  Even for the expanding circle map
$x\mapsto dx$, this quotient is not a group endomorphism and is not a fixed
Cartesian power.  The collision strata interact with periodic counting in a
particularly rigid way.

Write $\T=\mathbb R/\mathbb Z$ and let $X_k$ be the space of all nonempty
subsets of $\T$ of cardinality at most $k$.  For $d\geq2$, define
\[
 H_{d,k}:X_k\longrightarrow X_k,
 \qquad H_{d,k}(A)=\{dx\pmod1:x\in A\}.
\]
We give an exact unsigned periodic classification of this quotient, with a
standard finite-to-one entropy control.  Three conclusions carry the main
content.

\begin{enumerate}
\item We classify every fixed configuration of every iterate and derive a
  bivariate Euler product for the exact-cardinality strata.  Its rational
  collapse gives the closed formula in the abstract.
\item The total fixed count has an even--odd split, and the complete
  Artin--Mazur zeta function is an alternating product of $k$ linear factors
  (of $k+1$ factors when $k$ is odd and the factor at $d^0$ is included).
  M\"obius inversion gives the least-period orbit census and the asymptotic
  $d^{km}/m+O_{d,k}(d^{(k-1)m}/m)$.
\item The quotient map $\T^k\to X_k$ has uniformly finite fibers.  It
  therefore transfers the product entropy exactly, giving $k\log d$.  The
  two outermost zeta factors recover $d$ and $k$, so the parameter family is
  rigid under topological conjugacy.
\end{enumerate}

The first signal is visible without asymptotics.  At time $n$, put $Q=d^n$.
The counts for subsets of sizes $1,2,3,4$ are
\[
 Q-1,\qquad (Q-1)^2,\qquad (Q-1)(Q^2-Q+1),\qquad
 (Q-1)^2(Q^2+1).
\]
Their partial sums are $Q-1$, $Q(Q-1)$,
$(Q-1)(Q^2+1)$, and $Q(Q-1)(Q^2+1)$.  The alternating boundary term is the
early anomaly that leads to the parity theorem.

\section{Finite-subset spaces and the owner boundary}

For a compact metric space $Y$, its $k$th finite-subset space is
\[
 \exp_k(Y)=\{A\subseteq Y:1\leq\card A\leq k\},
\]
with the Hausdorff topology.  Equivalently, it has the quotient topology of
\[
 \pi_k:Y^k\longrightarrow \exp_k(Y),\qquad
 \pi_k(y_1,\ldots,y_k)=\{y_1,\ldots,y_k\}.
\]
We use $X_k=\exp_k(\T)$.  Tuffley determined the topology of these circle
spaces, their homotopy type, and the degree of maps induced from circle maps
\cite{Tuffley2002}.  In particular, $X_k$ has topological dimension $k$;
for $k\geq4$ it is generally a stratified compactum rather than a manifold.

The notation needs a firewall.  Some hyperspace papers call $\exp_k(Y)$ an
``$n$-fold symmetric product.''  The topological symmetric power
$\operatorname{SP}^k(Y)=Y^k/S_k$ remembers multiplicities, whereas
$\exp_k(Y)$ forgets them.  The two agree for $k\leq2$ but differ for
$k\geq3$: for example, $(a,a,b)$ and $(a,b,b)$ define different points in
$\operatorname{SP}^3(Y)$ and the same point in $\exp_3(Y)$.

General induced hyperspace dynamics is well developed.  Akin, Auslander,
and Nagar study the relation between a map and the induced map on compact
subsets \cite{AkinAuslanderNagar2017}.  Higuera and Illanes record, among
other facts, the finite-permutation mechanism behind periodic induced
subsets \cite{HigueraIllanes2011}.  G\'omez-Rueda, Illanes, and M\'endez
prove qualitative transfer statements for periodicity, recurrence, and
related properties on finite-subset spaces
\cite{GomezRuedaIllanesMendez2012}.  Fern\'andez, Good, and Puljiz analyze
admissible period sets in hyperspaces and symmetric products
\cite{FernandezGoodPuljiz2018}.  These results supply the qualitative owner
boundary; they do not give the unsigned exact-cardinality ledger computed
below.

For multiplication maps of the circle, Tan studies rotational finite subsets
and their rotation data \cite{Tan2024}.  Our count instead includes every
union of base cycles that fits below a cardinality cutoff, without imposing
one cyclic order type or a rotation number.

Finally, Artin--Mazur counting records cardinalities of fixed sets
\cite{ArtinMazur1965}.  Lefschetz zeta functions and symmetric-power
identities record fixed point indices; recent functorial treatments include
\cite{BlancoGomezEtAl2016,Crabb2025}.  We compare the two ledgers in
\cref{sec:entropy-control}, but do not identify them.  The residual theorem
package here is the circle-multiplication specialization and rational collapse
of binary cycle selection, its parity split, the resulting unsigned
Artin--Mazur factors, temporal census, and parameter rigidity.  The
finite-to-one entropy argument is a standard application of Bowen's factor
inequality, not a claimed family-specific mechanism.  A bounded search through 28 August
2026 did not locate that combined package.  This is not a priority proof,
and external novelty language remains on hold.

\section{Invariant finite sets as binary cycle selections}

Fix $d\geq2$ and $k\geq1$.  At time $n\geq1$, write
\[
 Q=d^n,
 \qquad g_Q=m_d^n=m_Q:\T\longrightarrow\T.
\]
For $\ell\geq1$, let $O_\ell(Q)$ denote the number of $g_Q$-orbits of
least period $\ell$.

\begin{lemma}[base orbit inventory]\label{lem:base}
For every $\ell\geq1$,
\begin{equation}\label{eq:base-orbits}
 \card{\Fix(g_Q^\ell)}=Q^\ell-1,
 \qquad
 O_\ell(Q)=\frac1\ell\sum_{e\mid\ell}
 \mob\!\left(\frac{\ell}{e}\right)(Q^e-1).
\end{equation}
Moreover, as a formal power series,
\begin{equation}\label{eq:base-zeta}
 \prod_{\ell\geq1}(1-u^\ell)^{-O_\ell(Q)}
 =\exp\!\left(\sum_{r\geq1}\frac{Q^r-1}{r}u^r\right)
 =\frac{1-u}{1-Qu}.
\end{equation}
\end{lemma}

\begin{proof}
The equation $g_Q^\ell(x)=x$ is
$(Q^\ell-1)x=0$ in $\T$.  Its solutions form the cyclic subgroup of order
$Q^\ell-1$.  If $R_\ell(Q)$ is the number of points of least $g_Q$-period
$\ell$, then
$Q^\ell-1=\sum_{e\mid\ell}R_e(Q)$.  M\"obius inversion gives
$R_\ell(Q)=\ell O_\ell(Q)$ and proves \eqref{eq:base-orbits}.

The infinite product in \eqref{eq:base-zeta} is well defined formally,
because only finitely many factors affect each coefficient.  Its logarithm
has coefficient
\[
 [u^r]\sum_{\ell\geq1}O_\ell(Q)
       \sum_{a\geq1}\frac{u^{a\ell}}a
 =\frac1r\sum_{\ell\mid r}\ell O_\ell(Q)
 =\frac{Q^r-1}{r}.
\]
Exponentiating gives the first equality.  Summing the two geometric
logarithms gives the rational expression.
\end{proof}

We denote the common formal series in \eqref{eq:base-zeta} by
$\zeta_{g_Q}(u)$.

The next lemma is the point where the induced dynamics becomes a binary
Euler transform.  The finite-permutation observation is general and appears
in the induced finite-subset literature; see, for example,
\cite[Lemma~4.7]{HigueraIllanes2011}.  We include the short proof to fix the
exact orbit inventory used below.

\begin{lemma}[orbit-union classification]\label{lem:union}
Let $A\subset\T$ be nonempty and finite.  Then $g_Q(A)=A$ if and only if
$A$ is a disjoint union of complete finite $g_Q$-orbits.  This decomposition
is unique.
\end{lemma}

\begin{proof}
If $g_Q(A)=A$, the restriction $g_Q|_A:A\to A$ is surjective.  A surjection
of a finite set is bijective, so $g_Q|_A$ is a permutation.  Its permutation
cycles are complete $g_Q$-orbits in $\T$, they are pairwise disjoint, and
their union is $A$.  The cycle decomposition of a permutation is unique.
Conversely, $g_Q$ maps every complete orbit onto itself, so it maps any
finite disjoint union of such orbits onto that union.
\end{proof}

For $j\geq1$, define the exact-cardinality fixed stratum
\[
 E_j(Q)=\card{\{A\in X_j:\card A=j,\ g_Q(A)=A\}}.
\]
The number is independent of the ambient cutoff once that cutoff is at least
$j$.  These strata are used only to partition fixed points; the exact-$j$
stratum is not generally forward invariant because collisions may reduce
cardinality.

\begin{theorem}[binary Euler quotient]\label{thm:euler}
As a formal power series in $u$,
\begin{align}
 \Phi_Q(u)
 &:=1+\sum_{j\geq1}E_j(Q)u^j \
 &=\prod_{\ell\geq1}(1+u^\ell)^{O_\ell(Q)}
   =\frac{\zeta_{g_Q}(u)}{\zeta_{g_Q}(u^2)}
   =\frac{1-Qu^2}{(1-Qu)(1+u)}.\label{eq:euler-quotient}
\end{align}
Consequently, for every $j\geq1$,
\begin{equation}\label{eq:exact-j}
 E_j(Q)=\frac{(Q-1)(Q^j-(-1)^j)}{Q+1}.
\end{equation}
\end{theorem}

\begin{proof}
By \cref{lem:union}, a fixed finite subset is specified by deciding, for
each $g_Q$-orbit, whether to include that orbit.  An orbit of length $\ell$
contributes either $1$ or $u^\ell$, and there are $O_\ell(Q)$ such orbits.
Multiplying these independent binary choices proves the first product.

Since $1+u^\ell=(1-u^{2\ell})/(1-u^\ell)$, \cref{lem:base} gives
\[
 \prod_{\ell\geq1}(1+u^\ell)^{O_\ell(Q)}
 =\frac{\prod_\ell(1-u^{2\ell})^{O_\ell(Q)}}
        {\prod_\ell(1-u^\ell)^{O_\ell(Q)}}
 =\frac{\zeta_{g_Q}(u)}{\zeta_{g_Q}(u^2)}.
\]
Substitution of \eqref{eq:base-zeta} cancels $1-u$ and produces the final
rational function in \eqref{eq:euler-quotient}.

For coefficient extraction, rewrite it as
\[
 \Phi_Q(u)=1+\frac{(Q-1)u}{(1-Qu)(1+u)}.
\]
The coefficient of $u^{j-1}$ in the reciprocal product is
\[
 \frac{Q}{Q+1}Q^{j-1}+\frac1{Q+1}(-1)^{j-1}
 =\frac{Q^j-(-1)^j}{Q+1}.
\]
Multiplication by $Q-1$ proves \eqref{eq:exact-j}.  The quotient is an
integer because $Q^j-(-1)^j$ is divisible by $Q+1$; integrality also follows
from the orbit-selection interpretation.
\end{proof}

\section{Parity-split fixed counts and zeta factors}

Put
\[
 A_k(Q)=\sum_{j=1}^k E_j(Q).
\]
Since $H_{d,k}^n(A)=g_Q(A)$ for $Q=d^n$, this is exactly
$\card{\Fix(H_{d,k}^n)}$.

\begin{theorem}[parity collapse]\label{thm:parity}
For $Q\geq2$ and $k\geq1$,
\begin{equation}\label{eq:total-fix}
 A_k(Q)=
 \begin{cases}
 \displaystyle \frac{Q(Q^k-1)}{Q+1},&k\ \text{even},\\[6pt]
 \displaystyle \frac{Q^{k+1}-1}{Q+1},&k\ \text{odd}.
 \end{cases}
\end{equation}
Equivalently, with
\[
 r_0(k)=\begin{cases}1,&k\ \text{even},\\0,&k\ \text{odd},\end{cases}
\]
one has the alternating polynomial
\begin{equation}\label{eq:alternating-polynomial}
 A_k(Q)=\sum_{r=r_0(k)}^k(-1)^{k-r}Q^r.
\end{equation}
\end{theorem}

\begin{proof}
Summing \eqref{eq:exact-j} gives
\[
 A_k(Q)=\frac{Q-1}{Q+1}
 \left(\frac{Q(Q^k-1)}{Q-1}-\sum_{j=1}^k(-1)^j\right).
\]
The alternating sum is zero for even $k$ and $-1$ for odd $k$.  This proves
\eqref{eq:total-fix}.  Polynomial division by $Q+1$ gives
\eqref{eq:alternating-polynomial}; its lower endpoint is $Q^1$ in the even
case and $Q^0$ in the odd case.
\end{proof}

Define the Artin--Mazur zeta function formally by
\[
 \zeta^{\AM}_{d,k}(z)=
 \exp\!\left(\sum_{n\geq1}
 \card{\Fix(H_{d,k}^n)}\frac{z^n}{n}\right).
\]

\begin{theorem}[alternating Artin--Mazur factors]\label{thm:am-zeta}
For $d\geq2$ and $k\geq1$,
\begin{equation}\label{eq:am-zeta}
 \boxed{\displaystyle
 \zeta^{\AM}_{d,k}(z)=
 \prod_{r=r_0(k)}^k
 (1-d^r z)^{(-1)^{k-r+1}}.}
\end{equation}
In particular,
\begin{center}
\begin{tabular}{c@{\qquad}c}
\toprule
$k$ & $\zeta^{\AM}_{d,k}(z)$\\
\midrule
$1$ & $\dfrac{1-z}{1-dz}$\\[5pt]
$2$ & $\dfrac{1-dz}{1-d^2z}$\\[5pt]
$3$ & $\dfrac{(1-z)(1-d^2z)}{(1-dz)(1-d^3z)}$\\[5pt]
$4$ & $\dfrac{(1-dz)(1-d^3z)}{(1-d^2z)(1-d^4z)}$\\
\bottomrule
\end{tabular}
\end{center}
\end{theorem}

\begin{proof}
Insert $Q=d^n$ into \eqref{eq:alternating-polynomial}.  Then
\begin{align*}
 \log\zeta^{\AM}_{d,k}(z)
 &=\sum_{r=r_0(k)}^k(-1)^{k-r}
   \sum_{n\geq1}\frac{(d^rz)^n}{n}\\
 &=-\sum_{r=r_0(k)}^k(-1)^{k-r}\log(1-d^rz).
\end{align*}
Exponentiation gives \eqref{eq:am-zeta}.  The displayed examples are direct
specializations.
\end{proof}

\begin{corollary}[parameter rigidity]\label{cor:rigidity}
Within the family $d\geq2$, $k\geq1$, the Artin--Mazur zeta function
determines the ordered pair $(d,k)$.  Hence $H_{d,k}$ and $H_{d',k'}$ can be
topologically conjugate only if $(d,k)=(d',k')$.
\end{corollary}

\begin{proof}
The factorization \eqref{eq:am-zeta} has no internal cancellation, because
$d^r$ is strictly increasing in $r$.  Its positive pole nearest the origin
is $d^{-k}$, and its positive zero nearest the origin is $d^{-(k-1)}$.  For
$k=1$ this zero is the factor at $d^0=1$.  The
ratio of the zero location to the pole location is $d$.  The pole then
recovers $k$.  Topological conjugacy bijects the fixed sets of corresponding
iterates and therefore preserves the Artin--Mazur zeta function.
\end{proof}

\section{Least-period orbits}

Let $P_{d,k}(m)$ be the number of $H_{d,k}$-orbits of least temporal period
$m$.

\begin{theorem}[exact temporal census]\label{thm:temporal}
For every $m\geq1$,
\begin{equation}\label{eq:temporal-census}
 P_{d,k}(m)=\frac1m\sum_{e\mid m}
 \mob\!\left(\frac me\right)A_k(d^e),
\end{equation}
where $A_k$ is given by either \eqref{eq:total-fix} or
\eqref{eq:alternating-polynomial}.  This expression is a nonnegative
integer.  For fixed $d\geq2$ and $k\geq2$,
\begin{equation}\label{eq:prime-orbit}
 P_{d,k}(m)=\frac{d^{km}}m
 +O_{d,k}\!\left(\frac{d^{(k-1)m}}m\right).
\end{equation}
\end{theorem}

\begin{proof}
Let $R_{d,k}(m)$ count points of least $H_{d,k}$-period $m$.  Every point
fixed by the $n$th iterate has a unique least period dividing $n$, so
\[
 A_k(d^n)=\sum_{m\mid n}R_{d,k}(m).
\]
M\"obius inversion gives the numerator of \eqref{eq:temporal-census}.
Every least-$m$ orbit contains exactly $m$ points, hence
$P_{d,k}(m)=R_{d,k}(m)/m$.  This proves both the formula and its
nonnegative integrality without an auxiliary congruence argument.

Equation \eqref{eq:alternating-polynomial} gives
\[
 A_k(d^m)=d^{km}+O_{d,k}(d^{(k-1)m}).
\]
Every proper divisor $e$ of $m$ satisfies $e\leq m/2$.  Also
$A_k(d^e)\leq C_{d,k}d^{ke}$, and therefore
\[
 \sum_{\substack{e\mid m\\e<m}}A_k(d^e)
 \leq C_{d,k}\sum_{e=1}^{\lfloor m/2\rfloor}d^{ke}
 =O_{d,k}(d^{km/2})
 =O_{d,k}(d^{(k-1)m}),
\]
where the last inequality uses $k\geq2$.  Substitution in
\eqref{eq:temporal-census} proves \eqref{eq:prime-orbit}.
\end{proof}

\section{Entropy and two owner controls}\label{sec:entropy-control}

The finite-subset quotient is not a product, but it has a product cover with
uniformly bounded fibers.  This section records a general entropy control,
not a residual novelty claim for the circle family; related entropy questions
for induced hyperspace maps are surveyed and developed in
\cite{KwietniakOprocha2007}.

\begin{theorem}[finite-to-one entropy equality]\label{thm:entropy}
For $d\geq2$ and $k\geq1$,
\begin{equation}\label{eq:entropy}
 h_{\mathrm{top}}(H_{d,k})=k\log d.
\end{equation}
\end{theorem}

\begin{proof}
Let $M_{d,k}=m_d^{\times k}$ on $\T^k$.  The quotient map $\pi_k$ satisfies
\[
 \pi_k\circ M_{d,k}=H_{d,k}\circ\pi_k.
\]
It is onto.  If $A\in X_k$ has $j$ elements, then $\pi_k^{-1}(A)$ consists
of the $k$-tuples in which every member of $A$ occurs at least once.  In
particular,
\[
 \card{\pi_k^{-1}(A)}\leq j^k\leq k^k,
\]
uniformly in $A$.

Entropy does not increase under a factor, so
$h_{\mathrm{top}}(H_{d,k})\leq h_{\mathrm{top}}(M_{d,k})$.  Bowen's factor
inequality bounds the reverse difference by the supremum of the Bowen
entropies of the fibers \cite{Bowen1971}.  Every fiber above has at most
$k^k$ points, independent of orbit length, and hence has Bowen entropy zero.
Thus the two entropies are equal.  This is the standard uniformly
finite-to-one consequence of Bowen's factor inequality.  The product formula and
$h_{\mathrm{top}}(m_d)=\log d$ give
$h_{\mathrm{top}}(M_{d,k})=k\log d$.
\end{proof}

\paragraph{Homological control.}
Tuffley's results \cite{Tuffley2002} give
\[
 X_k\simeq S^{2\lceil k/2\rceil-1},
 \qquad
 \deg(H_{d,k})=d^{\lceil k/2\rceil}
 \quad\text{on top nonzero homology}.
\]
Since this homology dimension is odd, the Lefschetz number of the $n$th
iterate is
\[
 L(H_{d,k}^n)=1-d^{n\lceil k/2\rceil},
\]
and the corresponding Lefschetz zeta function is
\[
 \zeta^{\Lef}_{d,k}(z)
 =\exp\!\left(\sum_{n\geq1}L(H_{d,k}^n)\frac{z^n}{n}\right)
 =\frac{1-d^{\lceil k/2\rceil}z}{1-z}.
\]
For $k\geq2$, this homological scale differs from the unsigned growth scale
$d^{kn}$ in \cref{thm:parity}.  Thus the Artin--Mazur factors in
\eqref{eq:am-zeta} are not a reformulation of the Lefschetz or fixed-index
formulas owned by the cited literature.

\paragraph{Multiplicity control.}
Periodic points of symmetric-product mappings have a direct earlier
literature, including Rallis \cite{Rallis1983}; the following calculation is
only a comparison control.
On the genuine symmetric power $\operatorname{SP}^k(\T)$, equality of a
multiset with its image first forces its finite support to map onto itself.
Thus $g_Q$ permutes the support, and the multiplicity is constant along every
support cycle.  A fixed multiset may consequently use a complete $g_Q$-orbit
with any nonnegative common multiplicity.
Therefore its mass generating function is
\[
 \prod_{\ell\geq1}(1-u^\ell)^{-O_\ell(Q)}
 =\frac{1-u}{1-Qu}.
\]
The coefficient at mass $k\geq1$ is $Q^{k-1}(Q-1)$.  Hence the induced map
on $\operatorname{SP}^k(\T)$ has
\[
 \zeta^{\AM}_{\operatorname{SP}^k(m_d)}(z)
 =\frac{1-d^{k-1}z}{1-d^kz}.
\]
This agrees with the finite-subset result at $k\leq2$, when multiplicity and
support encode the same quotient, but it lacks the lower alternating factors
for $k\geq3$.  The comparison isolates the mechanism: arbitrary orbit
multiplicity gives the ordinary Euler transform, whereas forgetting
multiplicity leaves a binary transform and creates the parity boundary.

\section{Exact controls, boundary cases, and scope}

The registered script implements two independent finite probes.  The first
computes the M\"obius orbit counts $O_\ell(Q)$ and multiplies the truncated
binary Euler factors with exact integer arithmetic.  For $2\leq Q\leq8$ and
$j,k\leq9$, it compares every coefficient and partial sum with
\eqref{eq:exact-j} and \eqref{eq:total-fix}.  A separate formal-factor probe
checks every logarithmic coefficient and the outer pole/zero recovery for
$2\leq d\leq7$, $1\leq k\leq9$, and iterates through 15.  It also checks the
temporal M\"obius census, including $k=1$, and its divisibility.

The second probe works with actual rational circle points.  For a chosen
$(Q,k)$ it takes
\[
 L=\operatorname{lcm}_{1\leq\ell\leq k}(Q^\ell-1),
\]
decomposes multiplication by $Q$ on $\mathbb Z/L\mathbb Z$ into literal
cycles, and counts cycle selections by cardinality.  Small cases additionally
enumerate every subset of $\mathbb Z/L\mathbb Z$ of size at most $k$ and test
$QA=A$ directly.  The two routes share neither coefficient extraction nor
subset enumeration.

The assumptions are material.  We exclude $d=1$, because the identity has
uncountably many fixed configurations and its unsigned Artin--Mazur zeta
function is not defined by finite cardinalities.  The case $k=1$ is included
as a base control, but $k\geq2$ is the genuinely higher-dimensional family.
We exclude the empty subset; adding it would add one fixed point to every
iterate and multiply \eqref{eq:am-zeta} by $(1-z)^{-1}$.  Exact-cardinality
coefficients $E_j(Q)$ are fixed-point strata, not zeta functions of invariant
exact-$j$ subsystems.

This manuscript is an internal theorem record.  Public posting, submission,
and absolute priority language remain on external hold pending a broader
specialist and database search.

\bibliographystyle{plainnat}
\bibliography{references}

\end{document}
